\subsection{Inserted-site scalar extraction}
\label{subsec:peps_inserted_site_scalar_extraction}

The final comparison of \cite[Section~3, proof of Theorem~3]{Molnar2018NormalPEPS}
turns two region proportionalities into a single scalar at one site. Fix a finite
region $R$ and a site $v\notin R$, and write $S=R\cup\{v\}$ for the
region with $v$ adjoined. A local configuration $\eta$ at $v$ assigns a virtual
index to every edge incident to $v$. The label on its $v$-incident $R$-boundary
edges reads $\eta$ on the $v$-incident edges that bound $R$, while the base
boundary label $\mu$ fixes the remaining $v$-incident edges. Under inserted-site
consistency, the $R$-boundary label and the values fixed by $\mu$ on the
remaining $v$-incident edges determine $\eta$ completely.

\begin{lemma}[Inserted-site consistency pins a local configuration to its $R$-boundary label]%
    \label{lem:peps_localConfig_eq_of_insertConsistent}
    \lean{TNLean.PEPS.localConfig_eq_of_insertConsistent}
    \leanok
    For $v\notin R$ and $\mu\in\partial(R\cup\{v\})$, any two local
    configurations $\eta,\eta'\in\mc V_v$ consistent with $\mu$ that induce the
    same $R$-boundary label $b(\mu,\eta)=b(\mu,\eta')$ satisfy $\eta=\eta'$.
\end{lemma}
\begin{proof}\leanok
    On each $e\in\operatorname{Inc}(v)$, if $e\in\partial R$ then
    $\eta_e=b(\mu,\eta)_e=b(\mu,\eta')_e=\eta'_e$, and if $e\notin\partial R$
    then $e\in\partial(R\cup\{v\})$ so consistency with $\mu$ gives
    $\eta_e=\mu_e=\eta'_e$.
\end{proof}

\begin{lemma}[Inserted-site scalar extraction from two region proportionalities]%
    \label{thm:peps_component_eq_of_regionProportional}
    \lean{TNLean.PEPS.component_eq_of_regionProportional}
    \leanok
    \uses{lem:peps_localConfig_eq_of_insertConsistent}
    Suppose $D_A=D_C>0$, $\inj(T_{C,R})$ holds, and $A_R=c_RC_R$ and
    $A_S=c_SC_S$ for $S=R\cup\{v\}$ ($v\notin R$) with $c_R\ne 0$. Then for
    every $\eta\in\mc V_v$ and $\sigma$,
    \begin{align}
        A_v(\eta,\sigma)
         & =\frac{c_S}{c_R}\,C_v(\eta,\sigma).
        \label{eq:peps_inserted_ratio}
    \end{align}
\end{lemma}
\begin{proof}\leanok
    With $b(\mu,\eta)\in\partial R$ the boundary label induced by $\mu\in\partial S$
    and $\eta\in\mc V_v$, and $m>0$ the interior bond multiplicity,
    \begin{align}
        m\,A_S(\mu,\sigma)
         & =
        \sum_{\eta\ \text{consistent with}\ \mu}
        A_v(\eta,\sigma(v))
        A_R(b(\mu,\eta),\sigma|_R),
        \label{eq:peps_insert_factorization}
    \end{align}
    and similarly for $C$. Substituting $A_S=c_SC_S$ and $A_R=c_RC_R$ into
    (\ref{eq:peps_insert_factorization}) and dividing by $m$ gives
    $\sum_{\eta}(c_RA_v(\eta,\sigma(v))-c_SC_v(\eta,\sigma(v)))C_R(b(\mu,\eta),\sigma|_R)=0$.
    Since $\eta\mapsto b(\mu,\eta)$ is injective by
    Lemma~\ref{lem:peps_localConfig_eq_of_insertConsistent} and $T_{C,R}$ is
    linearly independent, each coefficient vanishes, giving
    (\ref{eq:peps_inserted_ratio}).
\end{proof}

\begin{lemma}[Per-vertex gauge relation from two two-block proportionalities]%
    \label{cor:peps_component_eq_gaugeVertex_of_twoBlockProportional}
    \lean{TNLean.PEPS.component_eq_gaugeVertex_of_twoBlockProportional}
    \leanok
    \uses{thm:peps_component_eq_of_regionProportional, def:peps_regionTwoBlock, def:applyGauge, def:gaugeVertex}
    Suppose $D_A=D_B>0$, let $C=B^X$ reindexed along $D_A=D_B$, and assume
    $\inj(T_{C,R})$. If $B_{A,R}=c_RB_{C,R}$ and $B_{A,S}=c_SB_{C,S}$ for
    $S=R\cup\{v\}$ ($v\notin R$) with $c_R\ne 0$, then for every $\eta,\sigma$,
    \begin{align}
        A_v(\eta,\sigma)
         & =\frac{c_S}{c_R}\,C_v(\eta,\sigma).
        \label{eq:peps_vertex_gauge_ratio}
    \end{align}
\end{lemma}
\begin{proof}\leanok
    By Definition~\ref{def:peps_regionTwoBlock}, $B_{A,R}=c_RB_{C,R}$ and
    $B_{A,S}=c_SB_{C,S}$ mean $A_R=c_RC_R$ and $A_S=c_SC_S$, so
    Lemma~\ref{thm:peps_component_eq_of_regionProportional} yields
    (\ref{eq:peps_vertex_gauge_ratio}).
\end{proof}
